\section{引言}
部分子分布在高能散射过程中以部分子的语言刻画核子的结构。它们是给出强子-强子或轻子-强子碰撞实验物理预言的一个基本要素。尽管人们已经投入大量努力从各种实验数据中提取部分子分布~\cite{Alekhin:2012ig,Gao:2013xoa,Radescu:2010zz,CooperSarkar:2011aa,Martin:2009iq,Ball:2012cx}，但由于部分子分布的非微扰本质，从强相互作用的基本理论——量子色动力学（QCD）——出发计算它们是一项困难的任务。人们或许会问：能否从格点 QCD 出发计算部分子分布？格点 QCD 迄今为止仍是处理 QCD 中非微扰现象唯一可靠的框架。在场论的语言中，部分子分布由非局域的光锥关联函数给出：它既可以被视为无限动量极限下（在施加紫外（UV）截断之前）夸克和胶子的密度，也可以被视为有限核子动量下部分子的光前（light-front）关联~\cite{Collins:2011zzd}。这样的关联依赖于时间，因此无法直接在格点上计算。过去的尝试主要集中于计算分布的局域矩。然而，出于技术上的原因，在格点上计算部分子分布的高阶矩变得相当困难~\cite{negele}。

最近，本文作者之一提出了在欧氏格点上直接计算部分子物理的方案。在该方案中，人们计算的不再是光锥分布，而是一个相关的量，可以称之为准分布（quasi-distribution）~\cite{footnote}。对于非极化夸克密度，该准分布为~\cite{Ji:2013dva}
\begin{eqnarray}\label{qUnpolDef}
   \tilde q(x,\mu^2, P^z) &=& \int \frac{dz}{4\pi} e^{izk^z}  \langle P|\overline{\psi}(z)
   \gamma^z \\
   && \times \exp\left(-ig\int^{z}_0 dz' A^z(z') \right)\psi(0) |P\rangle \ ,  \nonumber
\end{eqnarray}
其中 $x= k^z/P^z$。上述量不含时间且是非局域的，对于任意满足 $P^z\ll 1/a$ 的动量（$a$ 为格点间距）均可在格点上模拟。然而，其结果并不是从实验数据中提取出的光锥分布 $q(x, \mu^2)$。为了从前者恢复后者，需要找到如下类型的匹配条件（非单态情形）
\begin{equation}
   \tilde q_{\rm NS}(x, \mu^2, P^z) = \int dy\, Z(x/y, \mu/P^z) q_{\rm NS}(y, \mu^2) + {\cal O} ((M/P^z)^2) \ ,
\end{equation}
它对大的 $P^z$ 成立，其中匹配因子 $Z$ 完全是微扰的。上述关系也可以被视为一个因子化定理：所有软发散在两边全部抵消，而 $\tilde q(x, \mu^2, P^z)$ 中的所有共线发散都与光锥分布中的相同。上述关系是否在微扰论的所有阶都成立尚待证明。不过作为第一步，本文将在单圈计算中对此进行检验。当然，准分布的选取并不唯一，人们可以定义不止一种与 $\tilde q(x, \mu^2, P^z)$ 具有类似性质的准分布。这里我们将关注最适合格点 QCD 计算的最简单类型。

本文考虑非单态夸克分布，并在单圈水平上证明上述因子化。在整个过程中我们得到了不含时间的准分布与光锥分布之间的匹配因子 $Z$。该结果将有助于以领头对数精度、从准分布的格点 QCD 模拟中恢复光锥分布。类型为 $(1/P^z)^{n}$（$n\ge 1$）的幂次压低修正在此被忽略，将在另文讨论。

本文其余部分安排如下：第 2 节给出非极化准夸克分布的单圈计算结果；基于此结果，我们在第 3 节提出一个在单圈水平成立的因子化定理，并提取准分布与光锥分布之间的匹配因子；第 4 节给出夸克螺旋度与横率分布的结果；最后在第 5 节作结论。
